\subsection{Why a round-one surrogate is not a shortcut}\label{sec:counterfactual}

Round-state statistics predict failure within the decoder that produced them \cite{kim2021logistic,paiva2022pickyfix}. We tested whether a surrogate trained on the reference variant's first-round counts transfers to a different schedule. The model was a logistic fit on the counts of correct and incorrect flips, trained on 46 keys and evaluated on 46 disjoint keys with 3,000 errors each. The two counts were the selected features; candidates were ranked by validation C0 log-loss, and the predicted RR is the predicted C0 DFR divided by the predicted intervention DFR. Within the arm that supplied the labels it ranked instances almost perfectly (AUC 0.995); that AUC is measured on the validation set used for model selection, not on a final held-out test. Across schedules the fit failed: for held-out offsets $+1$, $+2$ and $+3$ it predicted RR 0.79, 0.16 and 0.04, while the paired measurements are 3.46, 3.86 and 3.97. The counts are post-treatment variables: a low count of correct flips marks a hard instance under C0 but deferred flips under a raised threshold. This particular surrogate is therefore a diagnostic for the schedule that produced it, not a selector for another schedule. We do not claim that every round-one statistic fails as a predictor.
